%% 第 22 章：实谱定理
%% 本文件由 tools/md2tex.py 自动生成，请勿直接编辑。

\chapter{实谱定理}


\begin{jthm}{实谱定理}{r:23:1}
设 $V$ 是有限维实内积空间，$T:V\to V$。

下列条件等价：

\begin{jenum}
\item %
$T$ 是自伴算子；
\item %
$V$ 存在一组由 $T$ 的特征向量构成的规范正交基；
\item %
$T$ 在某组规范正交基下具有对角矩阵。
\end{jenum}

这里对角矩阵的全部对角元都是实特征值。
\end{jthm}

\section{为什么实谱定理要求自伴而非正规}

自伴算子一定正规。

但实正规算子不一定有实特征值。例如二维旋转算子保范且正规，却没有非零实特征向量，因此不能在实规范正交基下对角化。

自伴性比正规性更强，它排除了旋转型的二维耦合，只保留沿实特征方向的伸缩。

\begin{jprop}{一个二次算子可逆}{r:23:2}
若 $T$ 自伴，$a\in\mathbb R$，且 $b\in\mathbb R$ 非零，则：

\[
(T-aI)^2+b^2I
\]

可逆。
\end{jprop}

\begin{jproof}{证明思路}
$T-aI$ 仍然自伴。

对任意非零 $v$：

\[
\left\langle
\bigl((T-aI)^2+b^2I\bigr)v,v
\right\rangle
=
\|(T-aI)v\|^2+b^2\|v\|^2.
\]

右边严格大于零。

所以该算子的核只有零向量。在有限维空间中，单射等价于可逆，因此它可逆。
\end{jproof}

\section{实自伴算子一定有实特征值}


\begin{jproof}{证明思路}
先将 $T$ 复化，得到复内积空间上的算子 $T_{\mathbb C}$。

复数域上，$T_{\mathbb C}$ 一定有特征值：

\[
\lambda=a+ib.
\]

若：

\[
b\ne0,
\]

则复化后的算子：

\[
(T_{\mathbb C}-aI)^2+b^2I
\]

会将对应特征向量映为零。

但它正是前一个可逆算子的复化，因此不可能有非零核，产生矛盾。

所以：

\[
b=0.
\]

于是 $T$ 存在实特征值。
\end{jproof}

\begin{jprop}{自伴算子与不变子空间}{r:23:3}
设 $T$ 自伴，且 $U$ 是 $T$ 的不变子空间。

则：

\[
U^\perp
\]

也是 $T$ 的不变子空间。
\end{jprop}

\begin{jproof}{证明思路}
任取：

\[
w\in U^\perp,
\qquad
u\in U.
\]

因为 $U$ 对 $T$ 不变，所以：

\[
T(u)\in U.
\]

由自伴性：

\[
\langle T(w),u\rangle
=
\langle w,T(u)\rangle
=
0.
\]

因此：

\[
T(w)\in U^\perp.
\]
\end{jproof}

\begin{jprop}{限制算子仍然自伴}{r:23:4}
在上述条件下：

\[
T|_U:U\to U
\]

与：

\[
T|_{U^\perp}:U^\perp\to U^\perp
\]

都是自伴算子。
\end{jprop}

\begin{jproof}{证明思路}
自伴定义中的内积等式，对 $U$ 或 $U^\perp$ 内的向量仍然成立；而两者又各自不变，所以限制后的算子确实仍映回相应子空间。
\end{jproof}

\begin{jproof}{实谱定理的证明思路}
由前置引理，$T$ 存在一个单位特征向量 $e_1$。

于是：

\[
\operatorname{span}(e_1)
\]

是一维不变子空间。

其正交补：

\[
\operatorname{span}(e_1)^\perp
\]

也对 $T$ 不变，且限制算子仍自伴。

在这个维数减少一的正交补上重复同样过程。最终得到：

\[
(e_1,\ldots,e_n)
\]

这组规范正交特征基。
\end{jproof}
